\scanpage{230}
sections horizontales au voisinage de $s_0$, telles que
\[
a_0(s_0)=a(s_0),\quad b_0(s_0)=b(s_0).
\]
On a
\[
\mu(a,b)=\underbrace{\mu(a_0,b_0)}_{\deleted{\int_{a_0}^{b_0}\unsure{\varphi}}\;\varphi(b_0)-\varphi(a_0)}
+\mu(b_0,b)-\mu(a_0,a)
\]
donc
\[
\mu(a,b)(s_0)=\bigl[\varphi(b(s_0))-\varphi(a(s_0))\unsure{]|)}
+\left(\sum_i D_{is}\int_{b_0(s)}^{b(s)}\omega_{is}
-\sum_i D_{is}\int_{a_0(s)}^{a(s)}\omega_{is}\right)_{s=s_0}.
\]
Il faut calculer les deux termes de droite, c'est une histoire purement locale.

\deleted{Notons qu'on a, grâce à la connexion de $X/S$
\[
da\in\Gamma(S,a^*(\unsure{\mathcal T}_{X/S})\otimes\Omega^1_S)
\]
d'où en contractant avec
\[
a^{\unsure{*|\vee}}(\omega)\in\Gamma(S,\underbrace{a^*(\Omega^1_{X/S})}_{\text{dual de }a^{\unsure{*|\vee}}(\unsure{\mathcal T}_{X/S})})
\]
une section, que nous noterons $a^*(\omega)$
\[
a^*(\omega)\in\Gamma(S,\Omega^1_S)
\]
Cette \deleted{section} \inserted{forme} est encore \uline{fermée}, \inserted{$\omega$ l'étant} [N.B. pour le voir définir $a^{\unsure{*|\vee}}(\omega)$ pour tte forme (degré quelconque, fermée ou non) sur $X$, et noter que $a^*$ commute à la différentielle ext.)}
\margin{\unsure{donc une} \unsure{formule|forme} ....!}
\margin{229}
